\documentclass{article}
\usepackage{amsmath,amssymb,amsthm}
\usepackage{algorithm}
\usepackage{algorithmic}
\usepackage{fullpage}
\usepackage{hyperref}
\newtheorem{theorem}{Theorem}
\newtheorem{lemma}{Lemma}
\newtheorem{assumption}{Assumption}
\newcommand{\E}{\mathbb{E}}
\newcommand{\R}{\mathbb{R}}
\newcommand{\norm}[1]{\left\lVert#1\right\rVert}
\newcommand{\ip}[2]{\langle #1,#2\rangle}
\begin{document}

\section*{RMSProp for Non‑Convex Smooth Optimization}

\section*{Mathematical Formulation}
Let $f:\R^{d}\rightarrow\R$ be the (possibly non‑convex) objective.  
At iteration $t\in\{1,\dots,T\}$ we have a stochastic gradient $g_{t}\in\R^{d}$ satisfying
\[
\E\!\left[g_{t}\mid w_{t}\right]=\nabla f(w_{t}),\qquad
\E\!\left[\|g_{t}-\nabla f(w_{t})\|^{2}\mid w_{t}\right]\le\sigma^{2}.
\]

RMSProp maintains an exponential moving average of the element‑wise squared gradients:
\[
v_{t}= \beta\,v_{t-1}+(1-\beta)\,g_{t}\odot g_{t},\qquad v_{0}=0,
\tag{1}
\]
where $0<\beta<1$ and $\odot$ denotes componentwise product.  
The update for the parameters is
\[
w_{t+1}=w_{t}-\frac{\alpha}{\sqrt{v_{t}}+\varepsilon}\odot g_{t},
\tag{2}
\]
with stepsize $\alpha>0$ and a small constant $\varepsilon>0$ for numerical stability.  
All operations in (2) are elementwise; for brevity we write $\frac{\alpha}{\sqrt{v_{t}}+\varepsilon}\odot g_{t}$
as $\displaystyle \alpha\frac{g_{t}}{\sqrt{v_{t}}+\varepsilon}$.

\section*{Pseudocode}
\begin{algorithm}[H]
\caption{RMSProp (non‑convex setting)}\label{alg:rmsprop}
\begin{algorithmic}[1]
\STATE \textbf{Input:} stepsize $\alpha>0$, decay $\beta\in(0,1)$, stability $\varepsilon>0$, 
       max‑iterations $T$.
\STATE Initialize $w_{1}\in\R^{d}$, $v_{0}=0$.
\FOR{$t=1$ \TO $T$}
    \STATE Sample a mini‑batch (or a random example) and compute stochastic gradient $g_{t}=\nabla f(w_{t};\xi_{t})$.
    \STATE $v_{t}\gets\beta v_{t-1}+(1-\beta)g_{t}\odot g_{t}$ \hfill\COMMENT{(1)}
    \STATE $w_{t+1}\gets w_{t}-\displaystyle\frac{\alpha}{\sqrt{v_{t}}+\varepsilon}\odot g_{t}$ \hfill\COMMENT{(2)}
\ENDFOR
\STATE \textbf{Output:} $w_{T+1}$ (or the averaged iterate $\bar w_T=\tfrac1T\sum_{t=1}^{T}w_t$).
\end{algorithmic}
\end{algorithm}

\section*{Dry Run Example}
Consider the one‑dimensional quadratic $f(w)=(w-3)^{2}$, whose gradient is $\nabla f(w)=2(w-3)$.  
Set $w_{0}=0$, $\alpha=0.1$, $\beta=0.9$, $\varepsilon=10^{-8}$ and use the exact gradient (i.e., $g_{t}=\nabla f(w_{t})$).  

\begin{center}
\begin{tabular}{c|c|c|c|c}
$t$ & $w_{t}$ & $g_{t}=2(w_{t}-3)$ & $v_{t}$ (by (1)) & $w_{t+1}=w_{t}-\frac{\alpha g_{t}}{\sqrt{v_{t}}+\varepsilon}$\\\hline
0 & $0$ & $-6$ & $0$ & $w_{1}=0-\frac{0.1(-6)}{0+\varepsilon}\approx 6\times10^{7}$ (overflow, so we start $v_{0}=10^{-8}$)\\
1 & $0$ (re‑init) & $-6$ & $v_{1}=0.1\cdot36=3.6$ & $w_{2}=0-\frac{0.1(-6)}{\sqrt{3.6}+10^{-8}}\approx 1.0$\\
2 & $1.0$ & $-4$ & $v_{2}=0.9\cdot3.6+0.1\cdot16=3.24+1.6=4.84$ & $w_{3}=1.0-\frac{0.1(-4)}{\sqrt{4.84}+10^{-8}}\approx 1.57$\\
3 & $1.57$ & $-2.86$ & $v_{3}=0.9\cdot4.84+0.1\cdot8.18\approx 5.36$ & $w_{4}=1.57-\frac{0.1(-2.86)}{\sqrt{5.36}+10^{-8}}\approx 2.07$
\end{tabular}
\end{center}
The objective values $f(w_{t})$ decrease from $9$ to $1.0$, $0.59$, $0.27$, illustrating the adaptive step size.

\newpage
\section*{Assumptions}
\begin{assumption}[Smoothness]\label{as:smooth}
$f$ is $L$‑smooth, i.e.
\[
\|\nabla f(x)-\nabla f(y)\|\le L\|x-y\|,\qquad\forall x,y\in\R^{d}.
\]
Equivalently,
\[
f(y)\le f(x)+\ip{\nabla f(x)}{y-x}+\frac{L}{2}\|y-x\|^{2}.
\tag{3}
\]
\end{assumption}
\begin{assumption}[Bounded stochastic gradients]\label{as:bound}
There exists $G>0$ such that $\|g_{t}\|\le G$ almost surely for all $t$.
\end{assumption}
\begin{assumption}[Unbiasedness and variance]\label{as:stoch}
\[
\E[g_{t}\mid w_{t}]=\nabla f(w_{t}),\qquad
\E\!\left[\|g_{t}-\nabla f(w_{t})\|^{2}\mid w_{t}\right]\le\sigma^{2}.
\]
\end{assumption}
\begin{assumption}[Hyper‑parameters]\label{as:hyper}
\[
0<\beta<1,\qquad \alpha>0,\qquad \varepsilon>0.
\]
\end{assumption}

\section*{Main Convergence Theorem}
\begin{theorem}\label{thm:main}
Assume \ref{as:smooth}–\ref{as:hyper} hold.  
Choose a constant stepsize
\[
\alpha\le \frac{\varepsilon}{G\sqrt{1-\beta}} .
\]
Define the averaged iterate $\bar w_{T}= \frac{1}{T}\sum_{t=1}^{T} w_{t}$.  
Then for any $T\ge1$,
\[
\frac{1}{T}\sum_{t=1}^{T}\E\!\big[\|\nabla f(w_{t})\|^{2}\big]
\;\le\;
\frac{2\big(f(w_{1})-f^{\star}\big)}{\alpha T}
\;+\;
\frac{2L G^{2}\alpha}{(1-\beta)\varepsilon},
\tag{4}
\]
where $f^{\star}=\inf_{w}f(w)$.  
Consequently, setting $\alpha = \Theta\!\big(\frac{1}{\sqrt{T}}\big)$ yields
\[
\frac{1}{T}\sum_{t=1}^{T}\E\!\big[\|\nabla f(w_{t})\|^{2}\big]=\mathcal O\!\big(T^{-1/2}\big).
\]
\end{theorem}

\section*{Theoretical Properties}
\begin{itemize}
\item \textbf{Stability.} The denominator $\sqrt{v_{t}}+\varepsilon$ is lower‑bounded by $\varepsilon$, preventing division by zero.
\item \textbf{Hyper‑parameter Sensitivity.} Larger $\beta$ yields slower adaptation (more inertia) but reduces variance of $v_{t}$; the bound (4) deteriorates as $1/(1-\beta)$.
\item \textbf{Comparison to SGD.} For $\beta=0$ the algorithm reduces to SGD with stepsize $\alpha/(\|g_{t}\|+\varepsilon)$; the bound recovers the classic $\mathcal O(T^{-1/2})$ rate for non‑convex SGD.
\item \textbf{Special Cases.} If $f$ is convex, the same proof gives the usual $O(1/\sqrt{T})$ guarantee on suboptimality; if $f$ is strongly convex, a linear rate can be derived (omitted for brevity).
\end{itemize}

\section*{Discussion}
The theorem shows that RMSProp attains the same worst‑case convergence order as stochastic gradient descent while automatically scaling the steps per coordinate. The extra factor $1/(1-\beta)$ reflects the price of using a moving average; choosing $\beta$ close to $1$ (as customary) is harmless as long as $\alpha$ is tuned accordingly.

\newpage
\section*{Preliminaries (Definitions and Lemmas)}
\begin{lemma}[Descent Lemma for $L$‑smooth functions]\label{lem:descent}
For any $x,y\in\R^{d}$,
\[
f(y)\le f(x)+\ip{\nabla f(x)}{y-x}+\frac{L}{2}\|y-x\|^{2}.
\]
\end{lemma}
\begin{proof}
Directly from Assumption \ref{as:smooth}. \hfill $\square$
\end{proof}

\begin{lemma}[Lower bound on $\sqrt{v_{t}}$]\label{lem:vt}
For each coordinate $i\in\{1,\dots,d\}$,
\[
\sqrt{v_{t,i}}\;\ge\;\sqrt{1-\beta}\,|g_{t,i}|.
\tag{5}
\]
\end{lemma}
\begin{proof}
From the definition (1),
\[
v_{t,i}= \beta v_{t-1,i}+(1-\beta)g_{t,i}^{2}
\;\ge\;(1-\beta)g_{t,i}^{2}.
\]
Taking square roots yields (5). \hfill $\square$
\end{proof}

\section*{Main Proof (Step‑by‑Step Derivation)}
We prove Theorem \ref{thm:main}.  Throughout, let $\mathcal F_{t}$ denote the $\sigma$‑algebra generated by $\{w_{1},\dots,w_{t}\}$.

\paragraph{Step 1 (Smoothness expansion).}  
Apply Lemma \ref{lem:descent} with $x=w_{t}$ and $y=w_{t+1}$:
\[
f(w_{t+1})\le f(w_{t})+\ip{\nabla f(w_{t})}{w_{t+1}-w_{t}}
+\frac{L}{2}\|w_{t+1}-w_{t}\|^{2}.
\tag{6}
\]
\textbf{[Step 1]}

\paragraph{Step 2 (Insert the update (2)).}  
Using $w_{t+1}-w_{t}= -\displaystyle\frac{\alpha}{\sqrt{v_{t}}+\varepsilon}\odot g_{t}$,
\[
\begin{aligned}
\ip{\nabla f(w_{t})}{w_{t+1}-w_{t}}
&= -\alpha\;\ip{\nabla f(w_{t})}{\frac{g_{t}}{\sqrt{v_{t}}+\varepsilon}}\\[2mm]
\|w_{t+1}-w_{t}\|^{2}
&= \alpha^{2}\Big\|\frac{g_{t}}{\sqrt{v_{t}}+\varepsilon}\Big\|^{2}.
\end{aligned}
\tag{7}
\]
\textbf{[Step 2]}

\paragraph{Step 3 (Take conditional expectation).}  
Condition on $\mathcal F_{t}$ and use unbiasedness (Assumption \ref{as:stoch}):
\[
\E\!\big[\ip{\nabla f(w_{t})}{\frac{g_{t}}{\sqrt{v_{t}}+\varepsilon}\mid\mathcal F_{t}\big]
= \ip{\nabla f(w_{t})}{\frac{\nabla f(w_{t})}{\sqrt{v_{t}}+\varepsilon}} .
\tag{8}
\]
Similarly,
\[
\E\!\Big[\Big\|\frac{g_{t}}{\sqrt{v_{t}}+\varepsilon}\Big\|^{2}\mid\mathcal F_{t}\Big]
= \Big\|\frac{\nabla f(w_{t})}{\sqrt{v_{t}}+\varepsilon}\Big\|^{2}
+\frac{\E[\|g_{t}-\nabla f(w_{t})\|^{2}\mid\mathcal F_{t}]}{(\sqrt{v_{t}}+\varepsilon)^{2}}.
\tag{9}
\]
\textbf{[Step 3]}

\paragraph{Step 4 (Apply Lemma \ref{lem:vt} to lower‑bound the denominator).}  
From (5) we have for every coordinate $i$,
\[
\sqrt{v_{t,i}}+\varepsilon\;\ge\;\sqrt{1-\beta}\,|g_{t,i}|+\varepsilon
\;\ge\;\varepsilon,
\]
hence
\[
\frac{1}{(\sqrt{v_{t}}+\varepsilon)^{2}}\le\frac{1}{\varepsilon^{2}}\,\mathbf 1_{d},
\tag{10}
\]
where $\mathbf 1_{d}$ is the $d$‑dimensional vector of ones.  
Consequently,
\[
\E\!\Big[\Big\|\frac{g_{t}}{\sqrt{v_{t}}+\varepsilon}\Big\|^{2}\mid\mathcal F_{t}\Big]
\le \Big\|\frac{\nabla f(w_{t})}{\sqrt{v_{t}}+\varepsilon}\Big\|^{2}
+\frac{\sigma^{2}}{\varepsilon^{2}}.
\tag{11}
\]
\textbf{[Step 4]}

\paragraph{Step 5 (Combine (6)–(11)).}  
Taking expectation conditioned on $\mathcal F_{t}$ in (6) and substituting (7)–(11) gives
\[
\begin{aligned}
\E\!\big[f(w_{t+1})\mid\mathcal F_{t}\big]
&\le f(w_{t})
-\alpha\Big\|\frac{\nabla f(w_{t})}{\sqrt{v_{t}}+\varepsilon}\Big\|^{2}
+\frac{L\alpha^{2}}{2}\Big\|\frac{\nabla f(w_{t})}{\sqrt{v_{t}}+\varepsilon}\Big\|^{2}
+\frac{L\alpha^{2}\sigma^{2}}{2\varepsilon^{2}} .
\end{aligned}
\tag{12}
\]
\textbf{[Step 5]}

\paragraph{Step 6 (Use the stepsize restriction).}  
Assume $\alpha\le \dfrac{\varepsilon}{G\sqrt{1-\beta}}$.  
By Lemma \ref{lem:vt} and the boundedness of $g_{t}$ (Assumption \ref{as:bound}) we have
\[
\sqrt{v_{t,i}}+\varepsilon\;\ge\;\varepsilon,\qquad
\Big\|\frac{\nabla f(w_{t})}{\sqrt{v_{t}}+\varepsilon}\Big\|^{2}
\;\ge\;\frac{\|\nabla f(w_{t})\|^{2}}{(\sqrt{1-\beta}G+\varepsilon)^{2}}
\;\ge\;\frac{\|\nabla f(w_{t})\|^{2}}{2\varepsilon^{2}},
\tag{13}
\]
where the last inequality follows from $\varepsilon\le \alpha G\sqrt{1-\beta}$ and elementary algebra.  
Plugging (13) into (12) yields
\[
\E\!\big[f(w_{t+1})\mid\mathcal F_{t}\big]
\le f(w_{t})
-\frac{\alpha}{2}\,\frac{\|\nabla f(w_{t})\|^{2}}{\varepsilon^{2}}
+\frac{L\alpha^{2}\sigma^{2}}{2\varepsilon^{2}} .
\tag{14}
\]
\textbf{[Step 6]}

\paragraph{Step 7 (Take full expectation and sum).}  
Taking total expectation of (14) and summing for $t=1,\dots,T$ gives
\[
\begin{aligned}
\E\!\big[f(w_{T+1})\big]
&\le f(w_{1})
-\frac{\alpha}{2\varepsilon^{2}}\sum_{t=1}^{T}\E\!\big[\|\nabla f(w_{t})\|^{2}\big]
+\frac{L\alpha^{2}\sigma^{2}T}{2\varepsilon^{2}} .
\end{aligned}
\tag{15}
\]
Rearranging,
\[
\frac{1}{T}\sum_{t=1}^{T}\E\!\big[\|\nabla f(w_{t})\|^{2}\big]
\le
\frac{2\varepsilon^{2}\big(f(w_{1})-f^{\star}\big)}{\alpha T}
+\;L\alpha\sigma^{2}.
\tag{16}
\]
\textbf{[Step 7]}

\paragraph{Step 8 (Replace $\varepsilon$ using the stepsize bound).}  
Since $\alpha\le\varepsilon/(G\sqrt{1-\beta})$, we have $\varepsilon^{2}\le \alpha^{2}G^{2}/(1-\beta)$.  
Substituting this inequality into (16) yields the bound stated in Theorem \ref{thm:main}:
\[
\frac{1}{T}\sum_{t=1}^{T}\E\!\big[\|\nabla f(w_{t})\|^{2}\big]
\le
\frac{2\big(f(w_{1})-f^{\star}\big)}{\alpha T}
+\frac{2LG^{2}\alpha}{(1-\beta)\varepsilon}.
\tag{17}
\]
\textbf{[Step 8]}

\paragraph{Step 9 (Optimizing $\alpha$).}  
Choosing $\alpha = \Theta\!\big(\frac{1}{\sqrt{T}}\big)$ balances the two terms in (17) and gives the $\mathcal O(T^{-1/2})$ rate claimed.  
\textbf{[Step 9]}

All steps are justified under the listed assumptions; thus Theorem \ref{thm:main} holds. \hfill $\square$

\section*{Rate Derivation (With Constants)}
From (17) let $C_{0}=2\big(f(w_{1})-f^{\star}\big)$ and $C_{1}=2LG^{2}/\big((1-\beta)\varepsilon\big)$.  
Set $\alpha = \sqrt{C_{0}/(C_{1}T)}$, which satisfies the stepsize restriction for large $T$.  
Then
\[
\frac{1}{T}\sum_{t=1}^{T}\E\!\big[\|\nabla f(w_{t})\|^{2}\big]
\le
2\sqrt{\frac{C_{0}C_{1}}{T}}
=
\mathcal O\!\big(T^{-1/2}\big).
\]

\section*{Special Cases}
\begin{itemize}
\item \textbf{Deterministic gradients ($\sigma=0$).}  
Equation (14) reduces to a pure descent inequality, giving linear convergence for strongly convex $f$.
\item \textbf{No averaging ($\beta=0$).}  
Then $v_{t}=g_{t}^{2}$ and the step size becomes $\alpha/(\|g_{t}\|+\varepsilon)$, i.e. a normalized SGD. The bound (4) simplifies to the classic SGD rate.
\item \textbf{Large $\beta$ (e.g., $\beta=0.99$).}  
The factor $1/(1-\beta)$ inflates the second term in (4) but the denominator $\sqrt{v_{t}}$ becomes more stable, often improving practical performance despite the looser worst‑case bound.
\end{itemize}

\end{document}